\section{Introduction}

Place the two indices of the tensor
\(e_0\otimes e_0+e_1\otimes e_1\) at sites \(x\) and \(y\).
The indices carry the same binary label. This tensor acts as a
\emph{wire} because contracting its first index with a two-dimensional index of
another tensor moves that index from \(x\) to \(y\). Can an invariant
of placed tensors assign the wire itself its Euclidean length? We ask
for one assignment that works simultaneously for all pairs of sites
and respects local linear maps on all tensors.

Strassen's asymptotic spectrum supplies the algebraic rules for such an
invariant~\cite{Strassen1988}. A \emph{spectral point} on a fixed finite
set of two or more sites assigns a nonnegative number \(F(T)\) to each tensor.
Applying a linear map at each site is a \emph{restriction} and cannot
increase \(F\). Tensor products multiply its values, and direct sums
add them. Here a direct sum places the two tensors in separate blocks at
every site. The normalization is \(F(0)=0\) and \(F(1)=1\).
Taking a logarithm turns multiplication into addition. We call a
positive multiple of \(\log_2 F(T)\) a \emph{price}, with
\(\log_2 0=-\infty\). Independent wires
then have additive prices, and contracting wires produces a tensor whose price is at most the sum of
their prices.

For a finite nonempty set \(S\) of sites, write
\(\Delta_S(q)=\sum_{j=1}^q\bigotimes_{s\in S}e_j\) for the tensor
that shares one \(q\)-valued label among those sites. Let \(W(S)\)
denote the \emph{mean width} of their convex hull. It is the average
length of a one-dimensional projection over uniformly distributed
directions in the affine span, multiplied by a constant that gives
a segment its length. A singleton has width zero. Equivalently, for a standard Gaussian vector \(g\)
in the finite-dimensional span of \(S\),
\[
                   W(S)=\sqrt{2\pi}\,\mathbb E\max_{s\in S}\langle g,s\rangle .
\]
Translating all the sites adds a centered Gaussian to the maximum and
hence leaves its expectation unchanged.

\begin{theorem}[Euclidean prices]\label{thm:main}
Let \(k\) be any field and \(K\) any real Hilbert space. There is a
compatible assignment \(H(T;x)\) to every tensor over \(k\) and every
placement \(x=(x_1,\ldots,x_r)\) of its indices in \(K\) with the
following properties.
\begin{enumerate}[leftmargin=*,itemsep=.35ex]
\item A binary wire between \(u\) and \(v\) costs \(\|u-v\|\).
More generally, \(H(\Delta_S(q))=W(S)\log_2 q\) for every nonempty
finite \(S\subset K\) and every integer \(q\ge1\).
\item For each finite set \(S\subset K\) of positive diameter there is
a spectral point \(F_S\) such that \(H(T;x)=W(S)\log_2 F_S(T)\) for
every tensor placed in \(S\). In evaluating \(F_S\), group the
indices at each site and adjoin a one-dimensional factor at each unused
site. Compatibility means that enlarging \(S\) leaves \(H(T;x)\)
unchanged.
\item For each nonzero \(T\), the price depends on \(x\) only through
the labeled distances \(\|x_i-x_j\|\), and \(H(T;tx)=tH(T;x)\) for
\(t\ge0\).
\end{enumerate}
\end{theorem}

\begin{figure}[htbp]
\centering
\begin{tikzpicture}[font=\small, site/.style={circle,inner sep=2pt,
        fill=white,draw=ink,thick}, edge/.style={draw=sea,very thick}]
\begin{scope}
\coordinate (a) at (0,0); \coordinate (b) at (3.2,0); \coordinate (c) at (1.15,1.6);
\filldraw[fill=ink!6,draw=ink!55,dashed] (a)--(b)--(c)--cycle;
\foreach \v in {a,b,c} \node[site] at (\v) {\(i\)};
\node[below=7pt] at (1.6,0) {One common bit};
\end{scope}
\begin{scope}[xshift=6.0cm]
\coordinate (a) at (0,0); \coordinate (b) at (3.2,0); \coordinate (c) at (1.15,1.6);
\draw[edge] (a)--node[midway,above]{\(i\)}(b)--node[midway,above right]{\(j\)}(c)--node[midway,above left]{\(k\)}cycle;
\foreach \v in {a,b,c} \fill[ink] (\v) circle (2pt);
\node[below=7pt] at (1.6,0) {Three independent bits};
\end{scope}
\end{tikzpicture}
\caption{Two tensors on the same triangle of perimeter \(P\). On the
left, the tensor is supported on triples \((i,i,i)\) and costs \(P/2\).
On the right, each edge carries its own binary index. The tensor product
of the three wires costs \(P\).}
\label{fig:triangle}
\end{figure}

The shared-label formula is forced. Every compatible spectral system
realizing the wire distances on a whole Euclidean flat gives the same
value for \(\Delta_S(q)\) (Theorem~\ref{thm:flatlaw}). For planar sites
whose convex hull is a polygon this value per bit is half the hull's
perimeter. Figure~\ref{fig:triangle} illustrates why it matters which
labels are shared. The formula also controls tensors carrying a common
label alongside other data. If a tensor already carries that label on
\(S\), appending a copy at \(z\) increases its price by at most
\((W(S\cup\{z\})-W(S))\log_2 q\) (Proposition~\ref{lem:submod}).
In particular, a copy can be inserted anywhere on the segment joining
two existing copies at zero additional price. For arbitrary nonzero
\(T\), the chosen price is a seminorm on placements modulo common
translation. Applying a linear contraction to every site can only
decrease this price (Section~\ref{sec:hilbert}).

\subsection{Why a spectral point is needed}

Across a cut of the sites, a tensor can be viewed as a matrix, and its
matrix rank is a spectral point. A binary wire has rank two across a
cut separating its endpoints and rank one across every other cut.
Thus an average of log ranks measures how frequently the endpoints are
separated. Integrating over hyperplane cuts in Euclidean space recovers
the length of a wire. For each direction, the thresholds that split \(S\) form an interval
whose length is the width of the projection. Integrating these
lengths recovers the mean width for the shared label.

The difficulty is extending these evaluations to a spectral point on
\emph{all} tensors at the sites. For a small example, divide the four
corners of a square by its vertical and horizontal midlines. Let \(U\)
be a horizontal wire and \(V\) a vertical wire. Their matrix ranks
across the two cuts are \((2,1)\) and \((1,2)\), respectively. The
geometric mean of the ranks is \(\sqrt2\) for each tensor, but on their
direct sum it is \(\sqrt{3\cdot3}=3\). Additivity would give
\(2\sqrt2\). Even on wires, the metric obtained by averaging these two
cuts is a rescaled Hamming square and violates the four-point inequality
\eqref{eq:cur-quad} for spectral prices.

Direct sums matter for network constructions. If two tensors have the
same output spaces, their ordinary sum is a restriction of their direct
sum. The price of a sum of \(m\) alternatives on sites \(S\) is
therefore at most their largest price plus \(W(S)\log_2 m\).
Compatibility lets us use this bound even when the alternatives use
different auxiliary sites. For a single graph tensor, the tensor
product of its edge wires~\cite{cvzgraph}, each edge of dimension \(q\)
contributes its length times \(\log_2 q\). The wire formula gives the exact price of the graph tensor, and the
sum bound controls the price of alternative wirings.

\subsection{Relation to earlier work}

Strassen introduced the asymptotic spectrum to study transformations of
large tensor powers, motivated in particular by matrix multiplication
\cite{Strassen87,Strassen1988}. Comparison at all spectral points
characterizes asymptotic restriction. The difficult part in concrete
applications is to understand which evaluations can occur together.
Strassen developed support functionals~\cite{Strassen91}.

Christandl, Vrana and Zuiddam constructed quantum functionals yielding
spectral points over the complex field~\cite{CVZ}. The relation between
these constructions has since been clarified by Hirai and by Sakabe,
Do\u{g}an and Walter~\cite{hirai,sdw}. For three-party tensors, Alman,
Li and Pratt study edge support functionals over arbitrary
fields~\cite{almanlipratt}. For higher-order tensors, Botero and
coauthors study obstructions to universality and the possibility of
prescribing selected evaluations~\cite{botero}. Wigderson and Zuiddam
survey the duality theory and the interpolation of graph-tensor
evaluations towards a single-mode flattening rank~\cite{WZ}.

A triangle of independent wires represents matrix multiplication.
Bounds for its asymptotic rank constrain the three wire prices, and a
shared-leg copying argument in~\cite{nine} is related to our
incremental bound. Closure of price bounds under tensor degeneration
follows from the polynomial sublevel method
of~\cite{SpectralSublevels}.

The other papers in this collection use different aspects of this
framework. Shared labels occur in the sunflower argument~\cite{oaS},
and high-dimensional placements distinguish cuts of arithmetic
computations~\cite{oaH}. Spectral comparison and network compression
enter the circuit lower bounds~\cite{oaP}. Further applications concern
Fourier transforms, integer multiplication on tapes, and matrix
multiplication~\cite{oaF,oaT,oaM}.

\subsection{Proof strategy}

There are two questions. If a spectral space contains a Euclidean
flat, which tensor prices does its geometry force? And why must there
be such flats in arbitrarily large dimensions? A \emph{spectral flat}
of dimension \(d\) is a copy of \(\mathbb R^d\) with compatible
spectral prices whose binary wires cost their Euclidean lengths. The
existence statement that drives the proof is the following.
\begin{theorem}[Existence of spectral flats]\label{thm:flats}
Over every field there is a spectral flat of every positive finite
dimension.
\end{theorem}

Part~I derives the geometry from the algebra. The wire prices of any
spectral point satisfy the triangle inequality because two wires joined
at a common end supply a wire between their other ends. One can insert
a new site at any prescribed position along an interval
(Theorem~\ref{thm:extension}). Exterior algebra shows that when several indices move along their
intervals together, a tensor's price is convex in their common
parameter.
A shifted crossing identity then bounds the wire prices in a
quadrilateral and gives nonpositive curvature in the sense of
Alexandrov (Theorem~\ref{thm:geometry}, following the terminology of~\cite{BH}). Its two
flattenings build on crossing decompositions in
\cite{Anthropic,Khoruzhii}. For the underlying Clifford algebra
see~\cite{meinrenken}.

On a Euclidean flat, interpolation implies that adding a site inside
the convex hull of the others leaves the price of their shared bit
unchanged. A copying inequality and tensor restrictions determine the
price on segments and triangles. Approximating the support function of a polytope by
linear combinations of triangle support functions determines it on
every finite set (Section~\ref{sec:flatprice}). Finally, given flats
of unbounded dimension, concentration on high-dimensional frames makes
the prices of congruent placements agree in the limit. This yields one
system on the prescribed Hilbert space, including when its dimension
is uncountable (Section~\ref{sec:hilbert}).

\paragraph{Assume there is a largest flat.}
The proof of Theorem~\ref{thm:flats} occupies Parts~II and~III. Field
transfer reduces its failure to a finite field with a largest flat
dimension \(k\ge1\). A configuration approaching the corners of a
Euclidean \((k+1)\)-cube could be filled in, rescaled, and passed to a
limit to give a larger flat. The bound on dimension thus excludes
those corners uniformly. The rest of the argument uses this geometric
obstruction twice (Figure~\ref{fig:map}).

\begin{figure}[htbp]
\centering
\begin{tikzpicture}[>=Stealth,font=\small,
  box/.style={draw=ink,line width=.65pt,rounded corners=2pt,align=center,
              text width=4.8cm,inner sep=5pt},
  both/.style={box,text width=8.2cm},
  leftbox/.style={box,fill=ink!4},rightbox/.style={box,fill=sea!5,draw=sea},
  flowarrow/.style={->,line width=.65pt}]
\node[both] (cube) {A largest flat dimension excludes a larger cube\\
                     (Section~\ref{sec:prices})};
\node[leftbox] (port) at (-2.85,-1.65) {Bounded-degree port graphs\\
                     (Section~\ref{sec:ports})};
\node[rightbox] (probe) at (2.85,-1.65) {Local tests guide wire routing\\
                     (Section~\ref{sec:protocol})};
\node[leftbox] (comp) at (-2.85,-3.35) {Compress exact networks using tagged permutations\\
                     (Sections~\ref{sec:motifs}--\ref{sec:exact})};
\node[rightbox] (tag) at (2.85,-3.35) {Approximate the tagged gates\\with borrowed wires\\
                     (Section~\ref{sec:width})};
\node[both] (sim) at (0,-5.15) {Simulate many copies at smaller size per copy\\
                     (Section~\ref{sec:descent})};
\node[both] (diag) at (0,-6.75) {Distribute the compressed simulation's trace on the original grid
                     and derive a contradiction (Section~\ref{sec:diagonal})};
\draw[flowarrow] (cube.south -| port.north)--(port.north);
\draw[flowarrow,sea] (cube.south -| probe.north)--(probe.north);
\draw[flowarrow] (port)--(comp);
\draw[flowarrow,sea] (probe)--(tag);
\draw[flowarrow] (comp.south)--(sim.north -| comp.south);
\draw[flowarrow,sea] (tag.south)--(sim.north -| tag.south);
\draw[flowarrow] (sim)--(diag);
\end{tikzpicture}
\caption{The proof by contradiction. The left branch reduces network
size exactly if it may use gates for tagged permutations. The right
branch implements those gates with a controlled error. Compression
also reduces communication when the final argument distributes the
trace of the simulation.}
\label{fig:map}
\end{figure}

\paragraph{Compress exact networks.}
Under cube exclusion, bounded-degree graphs control the mean squared
distance between all pairs of vertices by the mean squared distance
along their edges, for every placement in a spectral space
(Section~\ref{sec:ports}). The construction uses nonlinear spectral
calculus~\cite{MendelNaor,MendelNaorBarycentric,MendelNaorHadamard}.
These \emph{port graphs} restore bounded degree during compression.
A second, fixed wire configuration lets us replace a network by a sum
of networks with fewer bits crossing between groups of vertices.
To assemble that sum into one exact computation, we must keep track of
which tagged records have been routed to which vertices. A
tagged-permutation gate tests that the addresses form a permutation
and that their payloads match. With these gates, the wire replacement
gives an exact reduction in size per copy
(Proposition~\ref{prop:oracle-descent}). Related methods of compressing
wire assignments occur in algorithms for satisfiability
\cite[Appendix~D]{falseeth} and in crossing decompositions~\cite{Anthropic}.

\paragraph{Implement the permutation gates.}
The second use of cube exclusion is local. Place \(k\) pairs of distant
wire endpoints along perpendicular axes at each tested location. If
nearby locations retain the prescribed prices for these wires, their
motion can be read from the changes in distances to the far ends.
A transverse segment in the limit would produce the excluded cube.
These tests range over possible local spectral prices. A failed test
allows a local wire replacement, while
successful tests identify a new link whose effect on subsequent prices
pays for its use. This gives an exact routing identity
(Theorem~\ref{thm:protocol}). Terms that spend too little on links between the layers have low
enough price to discard. Routing the surviving shorter lists recursively gives an approximation for
\(2^h\) records with \(h\)-bit payloads using
\(O(2^h h^{1+\varepsilon})\) trivalent vertices for every fixed
\(\varepsilon>0\) (Theorem~\ref{thm:cheaptag}).

The approximation has an algebraic meaning. Let \(I\) be the tensor
product of wires borrowed in the \(k\) coordinate directions. A
network intended to compute \(f\) while returning those wires produces
a tensor \(X\) in the same spaces as \(f\otimes I\). Its error is the
ordinary difference \(X-f\otimes I\). After fixing the data indices,
it suffices for the simulation that the error have price less than
\(H(I)\) by a specified margin in each prescribed flat placement of
the borrowed wires. The tagged-permutation approximation has this
bound even with its data indices exposed, as required for substitution. The
margin pays for subsequent sums. Once the margin exceeds the premium for adding two errors, a tensor
cannot approximate two distinct scalar multiples of \(I\), since their
difference has the full price \(H(I)\). Inserting these approximate gates into the exact
reduction and reusing borrowed wires along grid lines gives
Theorem~\ref{thm:collapse}. A bounded-degree computation of size \(V\)
with sufficiently few external bits then has an approximation of amortized
size at most \(V/(\log V)^A\) for arbitrarily large fixed \(A\).

\paragraph{Program the contradiction.}
Run many steps of a universal weighted computation on a grid. Let \(Z\)
be its exact tensor and \(\widehat Z\) the tensor obtained by substituting
the smaller simulations. The error estimate gives
\(H(Z-\widehat Z)<H(I)\). At this size, the original grid has enough
capacity to evaluate the accelerated contraction itself, once its trace
has been compressed. Each pair of external wire ends can be routed to
the step that needs its input while the other steps still have their
borrowed wires. We can consequently program the exact contraction to
return \(\widehat Z+I\). Its error is then \(I\), contradicting the
strict price bound. Section~\ref{sec:diagonal} carries out this
diagonal argument, which has a precedent in the use of faster
algorithms to prove circuit lower bounds~\cite{williams}.
